\section{Experto en estrategia I: escala, necesidades del usuario, tecnología y producto, capacidad organizacional}

El capítulo anterior trató el sistema de acción; este capítulo entra en el canal del experto en estrategia. Al hablar de la reunión estratégica de 2025 en Yanqi Hu (雁栖湖), Li Xiang (李想) sitúa la estrategia en un marco de cuatro variables: en el centro está la escala y, alrededor, las necesidades del usuario, la tecnología y el producto, y la capacidad organizacional. La escala no es un simple número de ingresos, sino algo que modifica la estructura de usuarios, la complejidad del producto, la forma de colaborar dentro de la organización y el panorama competitivo externo. Cuando una empresa pasa de ingresos de diez mil millones a cien mil millones, y después avanza hacia una escala aún mayor, tiene que volver a aprender las mejores prácticas del sector.

\begin{figure}[H]
\centering
\includegraphics[width=0.96\textwidth]{org-as-moe.png}
\caption{La organización también funciona como un MoE: el CEO asigna expertos, metodologías, talento joven y ciclos de retroalimentación. Diagrama conceptual propio, elaborado a partir de la conversación de 01:25:36--02:03:20.}
\end{figure}

\begin{knowledgebox}{Lectura de la figura: la organización no es una suma de personas, sino la asignación de expertos}
La tecnología, el producto, la estrategia, la contratación de recién egresados, la metodología y la retroalimentación forman en conjunto la capacidad organizacional. La organización se parece a un MoE porque las capacidades de los distintos expertos deben asignarse correctamente; si el enrutamiento es erróneo, por muchos expertos que haya, sus esfuerzos se anulan entre sí.
\end{knowledgebox}

\subsection{De Toyota, GM, Google y Huawei a una metodología propia}

El apartado anterior presentó la organización como un sistema de asignación de expertos; esta sección vuelve a cómo se forma ese sistema. En sus primeros años, Li Auto (理想) estudió el sistema de trabajo de Toyota, el proceso de investigación y desarrollo de GM y los OKR de Google; más adelante analizó también a Huawei y a Apple. Este aprendizaje no consistió en copiar, sino en buscar métodos transferibles para cada etapa de escala. En la etapa del Li ONE (理想 ONE) lo necesario era entregar; en la etapa de ingresos de cien mil millones se requieren densidad de talento, coordinación organizacional y capacidades de plataforma; en la etapa de dispositivos terminales de la era de la AGI se necesitan modelos, sistema operativo, hardware y ecosistema.

\begin{warningbox}{La metodología no debe tratarse como una religión}
Aprender de Toyota, Google, Huawei o Apple no significa copiar el lema organizacional de alguna de ellas. La metodología debe servir a la escala actual, a la etapa del producto y a las carencias de capacidad. Cuando cambia la escala de la empresa, también cambian las necesidades del usuario y la forma de la organización, y un método antiguo puede pasar de ser un activo a ser una carga.
\end{warningbox}

\subsection{Organizaciones pequeñas y alta densidad de talento}

Si la metodología determina la forma de colaborar, la siguiente pregunta es quién la ejecuta. Li Xiang menciona en varias ocasiones los equipos pequeños, la alta densidad de talento y la proporción de contrataciones de recién egresados. Un cambio posible en la era de la IA es que más organizaciones pequeñas obtengan mayor capacidad intelectual y fuerza anímica mediante modelos, herramientas y metodologías, en lugar de crecer linealmente en número de personas. El ejemplo de DeepSeek también se usa para ilustrarlo: el equipo no tiene que ser grande, pero sí debe tener alta densidad de talento, una dirección unificada y un ciclo de ingeniería sólido de principio a fin.

\begin{importantbox}{Los cuatro niveles de la capacidad organizacional}
El primer nivel es la densidad de talento: si las personas son lo bastante capaces. El segundo es la metodología: si las personas capaces colaboran con un mismo lenguaje. El tercero es la retroalimentación: si el negocio real y los usuarios corrigen continuamente el juicio. El cuarto es la energía: si el equipo está dispuesto a sostener tareas difíciles de manera continua.
\end{importantbox}

\subsection{Resumen de la sección}

El canal del experto en estrategia entiende a Li Auto dentro de los cambios de escala. Una empresa no avanza linealmente del punto A al punto B, sino que, cada vez que supera un escalón de escala, debe reconfigurar las necesidades del usuario, la tecnología y el producto, y la capacidad organizacional.
